\chapter{Learning decisions and designing their evidence}
\label{r3:decision-learning}
An estimated causal response is an input to a decision problem. The
bank must specify which decisions it can change, which resources they
consume and whose value is averaged. Selecting a policy also changes
the statistical problem: the selected rule is random, and its empirical
value benefits from the same noise that made it win.

\section{From heterogeneous effects to constrained policy value}
\label{r3:policy-regret}
Let $\tau(x)=\E[Y(1)-Y(0)\mid X=x]$ and let $c(x)$ be the
incremental action cost in the same financial units. The gain of a
binary rule $\pi$ over treating nobody is
\[
 G(\pi)=\E[\pi(X)\{\tau(X)-c(X)\}],\qquad
 R(\widehat\pi)=\sup_{\pi\in\Pi}G(\pi)-G(\widehat\pi).
\]
This regret is relative to the allowed class $\Pi$. If capacity requires
$\E[\pi(X)k(X)]\le B$, it restricts the feasible class and may change
the ranking of customers. A Lagrange relaxation uses net contrast
$\tau-c-\lambda k$; equality of the relaxation and the original
problem requires appropriate convexity or randomized policies. A
finite deterministic tree class can have a duality gap. Estimated
capacity and uncertain cost create further constraints and uncertainty.

Under the static identification assumptions, the AIPW pseudo-outcome
\[
 \Gamma(O)=m_1(X)-m_0(X)
   +\frac{A}{e(X)}\{Y-m_1(X)\}
   -\frac{1-A}{1-e(X)}\{Y-m_0(X)\}-c(X)
\]
satisfies $\E[\Gamma\mid X]=\tau(X)-c(X)$. Thus maximizing
$\mathbb P_n[\pi(X)\widehat\Gamma]$ is a weighted classification
problem: positive scores favor action, negative scores favor inaction,
with magnitude determining the cost of a mistake. Squared error for
$\tau$ values all covariate locations differently from this decision
loss. Accurate overall outcome prediction is further removed from it.

The following elementary bound makes estimation and optimization costs
explicit. Conditional on a separate nuisance-training sample, suppose
there are $M$ fixed candidate rules, independent evaluation observations,
and $|\widehat\Gamma|\le B_\Gamma$. Define the conditional population
score error
\[
 b_n=\sup_{\pi\in\Pi}
 \left|\E[\pi(X)\widehat\Gamma\mid\text{training}]-G(\pi)\right|.
\]
If the selected rule is within $\epsilon_{\rm opt}$ of the maximum
empirical score, then, with conditional probability at least $1-\alpha$,
\begin{equation}
 R(\widehat\pi)\le
 2B_\Gamma\sqrt{\frac{2\log(2M/\alpha)}{n}}
       +2b_n+\epsilon_{\rm opt}.
 \label{r3:finite-regret}
\end{equation}
For each candidate, its score lies in $[-B_\Gamma,B_\Gamma]$.
The bounded-variable exponential inequality and a union bound give
uniform deviation at most the square-root term without the leading
2. Add and subtract the empirical objectives for the optimal and
selected rules, use approximate maximization, and apply that deviation
and $b_n$ twice to obtain the result. The same argument applies to a
class chosen on the independent training data. It does not apply after
an unrestricted search on the evaluation outcomes with $M$ recorded
only as the number of final survivors.

With clipped nuisance estimates $\widehat e\in[\eta,1-\eta]$,
the conditional AIPW bias derived earlier implies
\[
 b_n\le\eta^{-1}\|\widehat e-e\|_2
       \bigl(\|\widehat m_1-m_1\|_2+\|\widehat m_0-m_0\|_2\bigr).
\]
This is a product error, not the generally larger absolute prediction
error of the individual pseudo-outcomes. Uniformity over the policy
class is the feature needed by the regret argument.

\citet[Assumptions 1--4 and Theorem 1]{policy2021} gives a much
sharper asymptotic analysis for cross-fitted scores and classes with
controlled VC dimension. Its nuisance conditions include uniform
consistency and expected squared errors of order $a(n)n^{-\zeta_m}$
and $a(n)n^{-\zeta_g}$, where $a(n)\to0$ and
$\zeta_m+\zeta_g\ge1$. The weighting representer and treatment
effect are bounded; outcome noise is uniformly sub-Gaussian with
conditional variance bounded away from zero. The allowed VC growth
$n^\beta$ has $0<\beta<1/2$ and
$\beta\le\min(\zeta_m,\zeta_g)$. Under the stated variance controls,
the regret has the familiar $\sqrt{\mathrm{VC}(\Pi)/n}$ scale,
with the score-variance and logarithmic factors in the theorem.
The source permits approximate maximization with its objective gap
$o(n^{-1/2})$. These conditions cannot be replaced by ordinary
ATE consistency. In particular, weak-IV weights violate its bounded
representer requirement, and arbitrary dependent panel rows do not
satisfy its sampling model.

\section{Evaluation after policy selection}
\label{r3:policy-selection}
One defensible design uses separate samples for nuisance learning,
policy selection and final evaluation. Conditional on the first two,
the selected rule is fixed for the last sample; its fixed-policy
gradient can be used subject to that sample's nuisance-rate and
dependence conditions. Cross-fitting can improve data use, but every
observation's roles must remain explicit. Training a policy on all
customers and then cross-fitting only the outcome regressions does
not make the policy independent of their outcomes.

For a prespecified finite class, simultaneous value intervals offer
another route. If $P(\forall\pi\in\Pi:G(\pi)\in I_\pi)\ge1-\alpha$,
the interval indexed by any data-selected member covers on that same
event. A valid joint multiplier or a conservative allocation of error
probabilities can construct such intervals under their respective
sampling assumptions. Repeatedly choosing subgroups, outcomes and
horizons enlarges the family beyond the final policy names. A lower
confidence bound on improvement is a possible deployment criterion;
it is a decision rule with a specified loss of power, not an unbiased
estimate of the selected maximum.

The same distinction applies to POMDP planning. Maximizing value over
an estimated transition model adds an optimization layer to parameter
and identification uncertainty. A small Bellman residual certifies
numerical accuracy inside that model; it does not certify the
out-of-sample value of its selected policy.

\section{Adaptive assignment and a time-uniform interval}
\label{r3:adaptive-inference}
Consider one bounded outcome per newly arriving customer, rescaled to
$[0,1]$, with no interference. Fix their potential outcomes
$(Y_i(0),Y_i(1))$. Before assignment, let the known randomization
probability $p_i\in[\epsilon,1-\epsilon]$ be measurable from past
observations and current pretreatment information. It may adapt to
earlier results. Outcomes are revealed in the specified order. The
inverse-assignment contrast
\[
 Z_i=\frac{A_iY_i}{p_i}-\frac{(1-A_i)Y_i}{1-p_i}
\]
has conditional mean $\tau_i=Y_i(1)-Y_i(0)$. Therefore
$S_t=\sum_{i\le t}(Z_i-\tau_i)$ is a martingale. This finite
randomization argument does not require identically distributed
customers. Its estimand is their running average effect
$\bar\tau_t=t^{-1}\sum_{i\le t}\tau_i$.

Conditional on the past, $Z_i$ is within
$[-1/(1-p_i),1/p_i]$, a range of width
$w_i=1/[p_i(1-p_i)]$. The bounded-variable exponential inequality
gives
\[
 \E[\exp\{\lambda(Z_i-\tau_i)\}\mid\mathcal F_{i-1}]
       \le\exp(\lambda^2 w_i^2/8),\qquad
 V_t=\sum_{i\le t}w_i^2/4.
\]
Here the filtration includes the current pretreatment information
before its randomization; it must exclude its outcome. Consequently
$\exp(\lambda S_t-\lambda^2V_t/2)$ is a nonnegative
supermartingale for each real $\lambda$. Mixing over
$\lambda\sim N(0,1/\rho)$ for fixed $\rho>0$ evaluates a Gaussian
integral and gives
\[
 M_t=\sqrt{\frac{\rho}{\rho+V_t}}
       \exp\left\{\frac{S_t^2}{2(\rho+V_t)}\right\}.
\]
The mixture is also a nonnegative supermartingale with $M_0=1$.
Stopping it at its first crossing of $1/\alpha$, bounded in time,
and using nonnegativity proves
$P(\sup_tM_t\ge1/\alpha)\le\alpha$ by passage to the limit.
Solving that crossing inequality supplies the confidence sequence
\begin{equation}
 \bar\tau_t\in\left[\bar Z_t-\frac{u_\alpha(V_t)}{t},
                      \bar Z_t+\frac{u_\alpha(V_t)}{t}\right],\quad
 u_\alpha(v)=\sqrt{(v+\rho)\log\frac{v+\rho}{\rho\alpha^2}},
 \quad\text{simultaneously for all }t,
 \label{r3:confidence-sequence}
\end{equation}
with probability at least $1-\alpha$. This reconstructs the normal
mixture principle in \citet[Definition 1, Lemma 2 and Appendix A.3,
Proposition 5]{howard2021}. Its conservative range bound is chosen
for transparent derivation; the source develops more efficient
variance-adaptive and iterated-logarithm boundaries under their
additional constructions. Fixed-sample normal intervals have no such
simultaneous guarantee merely because the assignments were randomized.

The problem persists even when every separate test has finite-sample
validity. Consider independent fair coin tosses. At each $t\le200$,
reject fairness only when the exact binomial lower or upper tail has
probability at most $.025$. Each fixed-time test has size at most
$.05$. An integer-count recursion over paths that have not previously
crossed a rejection boundary gives an overall rejection probability
of $.2837909361$ by $t=200$. The recursion extends each surviving
path by both possible tosses and removes first crossings, so this
number is an exact finite distribution calculation rounded for display.
It demonstrates repeated testing, without relying on a normal
approximation or a simulated rejection frequency.

Predictable outcome regression can augment $Z_i$ while retaining its
conditional mean, but the range or moment bound must be rederived.
Fitting that regression to future outcomes invalidates this predictable
argument. Estimated observational propensities likewise need bias
control inside the martingale construction. Optional stopping validity
does not establish validity after selecting any outcome, subgroup or
policy. For delayed defaults, revealing only quickly resolved loans
can select on outcomes; an arrival-order filtration, mature-cohort
design or justified delayed-observation model is required.

\section{Clusters and interference are part of the design}
\label{r3:interference}
A branch rollout or household pricing action can affect several
accounts. Let independent clusters be indexed by $c$, with account
outcomes $Y_{ci}(\mathbf a_c)$. Partial interference means another
cluster's assignment does not affect this outcome. An exposure mapping
further assumes
$Y_{ci}(\mathbf a_c)=Y_{ci}(a_{ci},s_{ci}(\mathbf a_{c,-i}))$,
where $s$ might be the fraction of neighbors offered the product.
This mapping is a restriction on spillovers, not a fact inferred from
account identifiers. Defining the intervention unit and interference
is a prerequisite to consistency \citep[Fine Point 1.1]{hernan2020}.

For a finite population of $N$ accounts and a known randomization
design, let $p_{ci}(a,s)>0$ be each account's exposure probability.
Then the design-based estimator
\[
 \widehat\mu(a,s)=\frac1N\sum_{c,i}
 \frac{\ind\{A_{ci}=a,S_{ci}=s\}Y_{ci}}{p_{ci}(a,s)}
\]
is unbiased for $N^{-1}\sum_{c,i}Y_{ci}(a,s)$: for each account,
the indicator's expectation cancels its exposure probability. This
project derivation says exactly which mean is recovered. Its variance
depends on joint exposure probabilities; accounts sharing a branch
are not independent replicates. Equal branch weights instead target
an average branch, which differs from the account mean when branch
size and effects covary.

If the entire branch is assigned together, exposures $(1,1)$ and
$(0,0)$ may occur while $(1,0)$ and $(0,1)$ never do. The rollout
identifies their bundled contrast but cannot separately identify
individual and spillover effects. More accounts within the same
branches do not create the missing variation. A two-stage design
that randomizes branch saturation and then accounts can supply
additional exposures at an explicit operational cost.

An exact randomization test also requires a stated null. Under the
sharp null $Y_{ci}(\mathbf a_c)$ is identical for every assignment,
all missing potential outcomes can be filled with observed values.
Recomputing a statistic over the known randomization distribution
then gives a valid tail probability by that distribution itself.
A zero average-effect null does not impute all potential outcomes;
the same exact-null proof does not apply to it. Spillovers across
branches or through shared market prices further invalidate the
partial-interference design and require a broader network or
equilibrium model.

The resulting evidence contract must distinguish a fixed-policy
interval, a uniform class bound, a time-uniform sequence, and a
randomization test. Their probability statements quantify over
different objects. A generic field named ``confidence'' cannot
preserve the guarantee needed by a decision maker.
